\documentclass[preprint,review,12pt]{elsarticle}
\usepackage{amssymb,amsmath}
\usepackage{booktabs}
\usepackage{multirow}
\usepackage{array}
\usepackage{makecell}
\usepackage[table]{xcolor}
\usepackage{graphicx}
\usepackage{url}
\usepackage{microtype}
\usepackage{placeins}
\usepackage{float}
\newcolumntype{L}[1]{>{\raggedright\arraybackslash}p{#1}}
\emergencystretch=3em

\DeclareUnicodeCharacter{2009}{\,}
\DeclareUnicodeCharacter{00F6}{\"o}
\DeclareUnicodeCharacter{015F}{\c{s}}
\DeclareUnicodeCharacter{00C7}{\c{C}}
\DeclareUnicodeCharacter{011F}{\u{g}}
\DeclareUnicodeCharacter{00E7}{\c{c}}

\journal{Knowledge-Based Systems}

\begin{document}

\begin{frontmatter}

\title{When does spatial attention transfer help a tiny leukemia
classifier? A controlled, multi-lens distillation study}

\author[its,telkom]{Tanzilal Mustaqim}
\author[its]{Chastine Fatichah\corref{cor1}}
\ead{chastine@its.ac.id}
\author[its]{Nanik Suciati}
\author[tok]{Takashi Obi}
\author[tok,kw]{Joong-Sun Lee}
\cortext[cor1]{Corresponding author}

\affiliation[its]{organization={Department of Informatics, Institut Teknologi Sepuluh Nopember},
            city={Surabaya}, postcode={60111}, country={Indonesia}}
\affiliation[telkom]{organization={Informatics Study Program, Telkom University, Surabaya Campus},
            city={Surabaya}, postcode={60231}, country={Indonesia}}
\affiliation[tok]{organization={Institute of Science Tokyo},
            city={Tokyo}, postcode={152-8550}, country={Japan}}
\affiliation[kw]{organization={Software Engineering Department, Catholic Kwandong University},
            city={Gangneung}, postcode={25601}, country={South Korea}}

\begin{abstract}
Deep networks now screen acute lymphoblastic leukemia (ALL) from
blood-smear images at close to expert accuracy, but a clinical tool must
also be small enough to run beside a microscope, transparent enough to
trust, and stable enough to deploy. Knowledge distillation compresses a
large teacher into a compact student, yet whether transferring
intermediate spatial attention on top of output logits buys anything in
practice stays unclear: prior work reports aggregate accuracy gains
without isolating the conditions that produce them. We do not propose a
new architecture. We run a controlled study that pins down when and how
spatial attention transfer helps a compact student. We distill a
SwinV2-Tiny teacher into a 4.21-million-parameter MobileNetV3-Large
student, 6.6$\times$ fewer parameters and 20.6$\times$ fewer FLOPs, and
compare Attention Transfer Knowledge Distillation (AT-KD) against Vanilla
Knowledge Distillation (VKD) under one shared teacher, student, and split
on three datasets, with stratified five-fold cross-validation, four
statistical tests, and Grad-CAM and t-SNE diagnostics. Three results
follow. The compact student matches the transformer teacher, and AT-KD
and VKD are statistically equivalent in accuracy; the four tests show
this parity is the expected ceiling once the discriminative signal
saturates, not a failure of attention transfer. The real value of
attention transfer is in optimization: it acts as a training-only
regularizer that more than halves validation loss on the spatially
complex ALL-IDB-1, speeds convergence, and shrinks fold-to-fold variance
at zero inference cost. This benefit is conditional on spatial
complexity, clear on multi-cell images and negligible on single-cell
crops, which yields a concrete rule for when attention transfer earns its
training cost. Grad-CAM and t-SNE confirm more morphology-aligned
activations and tighter feature clusters. We release a fully reproducible
protocol and deployment guidance for distillation in leukemia image
analysis.
\end{abstract}

\begin{highlights}
\item Spatial attention transfer matches vanilla KD in accuracy on saturated ALL benchmarks.
\item Four statistical tests confirm AT-KD--VKD accuracy parity as an expected ceiling effect.
\item Attention transfer halves ALL-IDB-1 validation loss at zero inference cost.
\item The benefit appears only on multi-cell images, giving a concrete deployment rule.
\item A 4.21M-parameter student matches its SwinV2-Tiny teacher with 20.6x fewer FLOPs.
\end{highlights}

\begin{keyword}
Knowledge distillation \sep attention transfer \sep lightweight deep learning \sep
leukemia classification \sep model interpretability \sep model calibration
\end{keyword}

\end{frontmatter}

\section{Introduction}

Acute lymphoblastic leukemia (ALL) is among the most common childhood
cancers, and survival depends on early detection and accurate subtyping.
Diagnosis starts at the microscope, where a specialist scans bone-marrow
or peripheral blood-smear images for blast cells and reads their
morphology. This manual review is slow, needs trained experts, and varies
between observers, and in many regions a shortage of specialists delays
it \cite{ref1,ref2,ref3}. Automated screening can give a fast first read
and flag suspicious cases for a human expert.

Deep learning has pushed leukemia image analysis close to expert
performance, with convolutional and transformer networks leading recent
benchmarks
\cite{ref3,ref4,ref5,ref18,ref19,ref20,ref21,ref22,ref23,ref24,ref25,ref26}.
Accuracy alone does not make a model deployable. The strongest networks
carry tens of millions of parameters and need a graphics processing unit
with ample memory, so they cannot run in real time on the low-power
hardware beside a microscope; a model that wins on paper can be unusable
at the point of care \cite{ref6,ref7,ref27,ref28,ref29,ref30,ref31}.
Knowledge distillation (KD) is the usual way to close this gap: a large
frozen teacher trains a small student to match its softened outputs, and
only the compact student ships \cite{ref8}. Vanilla KD (VKD) transfers
knowledge at the logit level alone and never constrains the intermediate
spatial cues that morphological diagnosis rests on, so a student can copy
the teacher's prediction while attending to the wrong region
\cite{ref9,ref10,ref11,ref32}. Attention Transfer Knowledge Distillation
(AT-KD) adds a loss that aligns the teacher and student spatial attention
maps \cite{ref12}, and attention-based distillation has been tried across
medical imaging
\cite{ref13,ref14,ref15,ref37,ref38,ref39,ref40,ref41,ref42,ref43}.

This activity leaves the central question open: does spatial attention
transfer actually help a lightweight leukemia classifier, and under what
conditions? Prior studies report aggregate accuracy gains but rarely
isolate when attention transfer helps; they say little about
optimization stability, convergence, calibration, or interpretability,
and they often rest on a single optimistic train-test split
\cite{ref16,ref17,ref33,ref34,ref35,ref36}. The basic deployment
question, whether attention transfer is worth its extra training cost and
when, has no clear answer.

We answer it with a controlled comparison rather than a new architecture.
We distill a SwinV2-Tiny transformer teacher into a lightweight
MobileNetV3-Large student and pit AT-KD against VKD under an identical
protocol on three public datasets of differing spatial character
(ALL-IDB-1, ALL-IDB-2, and a three-class Subtype ALL set), with
stratified five-fold cross-validation, four statistical tests, and
Grad-CAM and t-SNE diagnostics. The headline finding is that attention
transfer does not raise accuracy over vanilla distillation, and that this
parity is the expected outcome at this scale; the payoff of attention
transfer lies instead in optimization and interpretability, and it tracks
the spatial complexity of the input.

This study makes four contributions.
\begin{itemize}
\item \textbf{A deployment decision rule.} Attention transfer earns its
extra training cost \emph{only} when the imagery carries multi-cell
spatial context: the gain is clear on high-resolution ALL-IDB-1 and
three-class Subtype ALL and vanishes on single-cell ALL-IDB-2. This
replaces a vague ``attention helps'' with a testable rule that prior work
does not provide.
\item \textbf{Attention transfer as a zero-cost optimization
regularizer.} We identify and mechanistically explain spatial attention
transfer as a training-only regularizer that more than halves validation
loss on ALL-IDB-1, speeds convergence, and cuts fold-to-fold variance at
no inference cost, and we show why this gain does not surface as accuracy
on saturated data.
\item \textbf{A controlled, multi-lens protocol.} Under one shared
teacher, student, and split we combine stratified five-fold
cross-validation, four statistical tests (Friedman, paired $t$-test,
McNemar, and effect sizes), and Grad-CAM and t-SNE diagnostics,
establishing AT-KD--VKD accuracy parity as a tested, expected result
rather than an unexamined claim.
\item \textbf{A deployable, interpretable student.} The distilled
4.21-million-parameter MobileNetV3-Large matches its SwinV2-Tiny teacher
at 6.6$\times$ fewer parameters and 20.6$\times$ fewer FLOPs, with more
morphology-aligned activations and tighter feature clusters.
\end{itemize}

\section{Related work}
Detecting ALL from microscopic blood images turns on subtle
morphological differences between subtypes, and getting the subtype right
drives treatment selection: a misclassification can lead to inappropriate
therapy and serious complications, especially in pediatric patients
\cite{ref1,ref6,ref16,ref17}. Computational approaches fall into three
groups, reviewed below: conventional feature engineering with classical
machine learning, transfer learning with pretrained deep networks, and
knowledge distillation for model compression. The recurring gap is that
prior work optimizes aggregate accuracy while leaving the conditions of
benefit, the interpretability, and the evaluation rigor of
attention-based distillation largely untested.

\subsection{Conventional and CNN-based approaches for ALL diagnosis}

Early computational methods relied on hand-crafted features such as the
nuclear-cytoplasmic ratio, area, compactness, roundness, and texture,
typically extracted by k-means clustering or watershed segmentation and
then classified with support vector machines, k-nearest neighbors, Naive
Bayes, or Random Forest \cite{ref18,ref19}. These pipelines reached
moderate accuracy but depended on expert-designed features and heavy
preprocessing, and they stayed sensitive to inter-observer variability,
dataset differences, and staining variation \cite{ref20,ref21}.

Basic convolutional neural networks (CNNs) such as AlexNet and VGGNet
reduced the dependence on manual feature extraction and improved
performance on benchmarks such as ALL-IDB \cite{ref4}. Their
effectiveness remained limited by the scarcity of labeled data, however,
leaving them prone to overfitting and sensitive to image resizing, which
can distort the fine cellular detail that leukemia diagnosis depends on
\cite{ref18,ref22}.

\subsection{Transfer learning for ALL subtype classification}

Transfer learning has since improved ALL classification and subtype
differentiation. Architectures such as VGG16, VGG19, ResNet50,
InceptionV3, DenseNet-121, and MobileNetV2 \cite{ref4,ref21,ref22} adapt
feature hierarchies learned on large-scale datasets to leukemia images
through fine-tuning \cite{ref23,ref24}, and recent work pairs lightweight
backbones with attention mechanisms or metaheuristic optimization to
raise accuracy on ALL benchmarks \cite{ref25,ref26}.

These models remain computationally demanding, often requiring high-end
GPUs and longer inference time, which limits their use in real-time and
point-of-care settings \cite{ref6}. Effective fine-tuning also depends on
large annotated datasets that are hard to obtain in medical imaging
because of privacy regulation, annotation cost, inter-observer
variability, and the scarcity of samples for rare subtypes
\cite{ref27,ref28,ref29,ref30}, and the limited interpretability of deep
models can erode clinician trust and slow clinical adoption \cite{ref31}.

\subsection{Knowledge distillation for model compression}

Knowledge distillation (KD) offers a practical route to model
compression \cite{ref8,ref32}: a compact student is trained to imitate
the temperature-scaled probability distribution of a larger teacher,
usually through Kullback--Leibler divergence \cite{ref8}. It has been
applied to histopathological and general medical image classification
\cite{ref33,ref34}, and in leukemia to compress teacher CNNs into compact
students \cite{ref35}, although response-based distillation transfers only
output-level knowledge and discards spatial reasoning.

Because vanilla KD operates at the logit level, it ignores intermediate
feature representations, spatial attention maps, and relational structure
\cite{ref9}. This matters for ALL subtyping, where the spatial
relationships among nuclei, cytoplasm, chromatin patterns, and blast-cell
clusters are the key discriminative cues in blood-smear images and are
poorly preserved by response-based distillation \cite{ref34}. A student
can then reproduce the teacher's prediction without acquiring the spatial
reasoning behind it, which weakens generalization on spatially complex
images.

\subsection{Attention mechanisms in medical imaging}

Spatial attention helps distinguish subtle morphological patterns among
ALL subtypes, including variation in the nuclear-cytoplasmic ratio,
chromatin texture, and blast-cell clustering \cite{ref2,ref36}.
Attention-based CNNs and regionally focused models improve ALL
discrimination by emphasizing nuclei and blast regions
\cite{ref37,ref38,ref39}, and across several medical-imaging tasks models
with stronger spatial awareness outperform their non-attention
counterparts \cite{ref40,ref41}.

Attention Transfer Knowledge Distillation (AT-KD) extends KD by aligning
teacher and student attention maps from intermediate layers \cite{ref12},
transferring spatial feature hierarchies alongside output-level knowledge
\cite{ref34}. It has been used for whole-slide image analysis, multi-task
tumor grading, skin lesion analysis, and retinal disease grading
\cite{ref13,ref14,ref15,ref42}, but its behavior in ALL classification
remains poorly characterized, particularly with respect to spatial
complexity, convergence, and interpretability.

\subsection{Research gap}

Across this literature, attention transfer is treated as a general
accuracy booster, and its specific role in lightweight ALL classification
has not been examined in detail. Three gaps stand out.

First, few studies connect the benefit of attention transfer to the
spatial characteristics of the input. Since ALL diagnosis relies on
morphological cues spread across nuclear structure, cytoplasm, chromatin
texture, and blast-cell distribution, attention transfer should help more
on high-resolution multi-cell images than on small single-cell crops, yet
this condition has not been isolated.

Second, deployment-related behavior is underexplored: most studies report
accuracy but say little about validation loss, convergence, optimization
stability, or interpretability, even though a deployable model must be
accurate, efficient, stable, and explainable at once.

Third, the frequent reliance on a single train-test split makes results
sensitive to how the data are partitioned and can mask whether a reported
gain is a consistent effect or fold-specific noise.

We address these gaps by characterizing AT-KD for deployment-oriented ALL
classification across datasets of differing spatial complexity, examining
convergence, validation loss, and visual interpretability rather than
accuracy alone. Stratified five-fold cross-validation with a multi-lens
statistical evaluation separates consistent behavior from cross-fold
variation. Unlike segmentation-supervised pipelines that need pixel-level
masks \cite{ref43}, our framework derives spatial supervision directly
from teacher activations, transferring attention without segmentation
annotation and staying practical where expert labeling is costly.

\section{Method}
Figures~\ref{fig:1}--\ref{fig:3} give an overview of the two distillation
schemes we compare and of the three benchmark datasets, ahead of the
formal description below.

\begin{figure}[tbp]
\centering
\includegraphics[width=0.85\linewidth]{Figures/fig1.png}
\caption{Vanilla Knowledge Distillation (VKD) architecture.}
\label{fig:1}
\end{figure}
\begin{figure}[tbp]
\centering
\includegraphics[width=0.85\linewidth]{Figures/fig2.png}
\caption{Proposed Attention Transfer Knowledge Distillation (AT-KD) architecture.}
\label{fig:2}
\end{figure}
\begin{figure}[tbp]
\centering
\includegraphics[width=0.7\linewidth]{Figures/fig3.png}
\caption{Representative samples from the three benchmark datasets.}
\label{fig:3}
\end{figure}

\subsection{Vanilla and attention-transfer knowledge distillation}

In vanilla knowledge distillation, a frozen high-capacity teacher
produces softened class probabilities that, together with the
ground-truth labels, supervise a compact student through a
temperature-scaled Kullback--Leibler divergence. The soft-target
probability for class $i$ follows \eqref{eq:1}, where $z_i$ is the logit
and $T$ the temperature; the total objective combines cross-entropy on
hard labels with the distillation term in \eqref{eq:2}, and the gradient
of the soft-target loss scales by $1/T^{2}$ as in \eqref{eq:3}.

\begin{equation}\label{eq:1}
q_{i} = \frac{\exp(z_{i}/T)}{\sum_{j}^{}{\exp(}z_{j}/T)}
\end{equation}
\begin{equation}\label{eq:2}
L_{\text{total}} = (1 - \alpha) \cdot L_{\text{CE}} + \alpha \cdot L_{\text{KL}}
\end{equation}
\begin{equation}\label{eq:3}
\frac{\partial C}{\partial z_{i}} = \frac{1}{T^{2}}(q_{i} - p_{i})
\end{equation}

Attention Transfer Knowledge Distillation augments this objective by
aligning the spatial attention maps of teacher and student. From a
feature tensor we form a spatial attention map by summing the squared
channel responses and normalizing, as in \eqref{eq:4} and \eqref{eq:10};
the activation-based and gradient-based attention-transfer losses appear
in \eqref{eq:5}--\eqref{eq:7}. The complete AT-KD objective adds the
morphology-aware attention-transfer term to the logit distillation and
cross-entropy terms, with the weights in \eqref{eq:8} and \eqref{eq:9}.
The attention-transfer terms act only during training, so the deployed
student is identical in size and inference cost to a vanilla-distilled
student.

\begin{equation}\label{eq:4}
F_{\text{sum}}^{p}(A) = \sum_{i = 1}^{C} \mid A_{i} \mid^{p},p > 1
\end{equation}
\begin{equation}\label{eq:5}
\begin{aligned}
L_{\text{AT}}^{\text{act}} &= L_{\text{CE}}(W_{S}, x) \\
&\quad + \frac{\beta}{2}\sum_{j \in I}\left\| \frac{Q_{S}^{j}}{\left\| Q_{S}^{j} \right\|_{2}} - \frac{Q_{T}^{j}}{\left\| Q_{T}^{j} \right\|_{2}} \right\|_{2}^{2}
\end{aligned}
\end{equation}
\begin{equation}\label{eq:6}
\begin{aligned}
L_{\text{AT}}^{\text{grad}} &= L_{\text{CE}}(W_{S}, x) + \frac{\beta}{2}\left\| J_{S} - J_{T} \right\|_{2}^{2}, \\
J_{S} &= \frac{\partial L_{\text{CE}}(W_{S}, x)}{\partial x}, \quad J_{T} = \frac{\partial L_{\text{CE}}(W_{T}, x)}{\partial x}
\end{aligned}
\end{equation}
\begin{equation}\label{eq:7}
\mathcal{L}_{\text{sym}} = \mathcal{L}_{\text{CE}}(W,x) + \frac{\beta}{2}\left\| J - \text{flip}(J) \right\|_{2}^{2}
\end{equation}
\begin{equation}\label{eq:8}
\begin{aligned}
\mathcal{L}_{\text{total}} &= \alpha_{\text{CE}}\mathcal{L}_{\text{CE}} + \beta_{\text{KL}}T^{2}\mathcal{L}_{\text{KL}}, \\
&\quad \alpha_{\text{CE}} = \beta_{\text{KL}} = 0.5
\end{aligned}
\end{equation}
\begin{equation}\label{eq:9}
\begin{aligned}
\mathcal{L}_{\text{total}} &= \alpha_{\text{CE}}\mathcal{L}_{\text{CE}} + \beta_{\text{KL}}T^{2}\mathcal{L}_{\text{KL}} + \gamma\mathcal{L}_{\text{AT}}, \\
&\quad \alpha_{\text{CE}} = 1.0,\ \beta_{\text{KL}} = 0.5,\ \gamma = 50.0
\end{aligned}
\end{equation}
\begin{equation}\label{eq:10}
\begin{aligned}
\mathcal{A}(F) &= \frac{\mathbf{v}}{\left\| \mathbf{v} \right\|_{2}}, \\
\mathbf{v} &= \text{vec}\!\left( \frac{1}{C}\sum_{c=1}^{C}\left| F_{c} \right|^{2} \right)
\end{aligned}
\end{equation}

Algorithm~\ref{alg:training} summarizes the full VKD and AT-KD procedure.
We first train the teacher, then freeze it and distill the student under
either the vanilla or the attention-transfer objective, selecting the
best checkpoint of each fold by macro-F1.

\begin{figure}[tbp]
\centering
\rule{\columnwidth}{0.9pt}\\[1pt]
{\footnotesize\textbf{Algorithm 1}\quad VKD and AT-KD Training Procedure}\\[2pt]
\rule{\columnwidth}{0.4pt}\\[1pt]
\footnotesize
\setlength{\tabcolsep}{3pt}
\renewcommand{\arraystretch}{1.15}
\begin{tabular}{@{}r@{\ \ }p{0.86\columnwidth}@{}}
1  & \textbf{Input:} Training set $D_{\text{train}}$, validation set $D_{\text{val}}$, teacher configuration $C_{t}$, student configuration $C_{s}$, KD method $M \in \{\text{Vanilla},\ \text{Attention Transfer}\}$ \\
2  & \textbf{Output:} Best student model $S_{\text{final}}$ \\
3  & \textbf{Initialize} teacher model $T$ using $C_{t}$ \\
4  & \textbf{for} $\text{epoch} = 1$ to $E_{t}$ \textbf{do} \\
5  & \hspace{1.5em}\textbf{Train} $T$ using cross-entropy loss on $D_{\text{train}}$ \\
6  & \hspace{1.5em}\textbf{Validate} $T$ on $D_{\text{val}}$ \\
7  & \textbf{end for} \\
8  & \textbf{Freeze} $T$ and set $T$ to evaluation mode \\
9  & \textbf{Initialize} student model $S$ using $C_{s}$ \\
10 & \textbf{Set} $\text{best\_score} \leftarrow 0$ \\
11 & \textbf{for} $\text{epoch} = 1$ to $E_{s}$ \textbf{do} \\
12 & \hspace{1.5em}\textbf{for} each mini-batch $(x, y)$ in $D_{\text{train}}$ \textbf{do} \\
13 & \hspace{3em}\textbf{Compute} student logits $z_{S} = S(x)$ \\
14 & \hspace{3em}\textbf{Compute} teacher logits $z_{T} = T(x)$ without gradient \\
15 & \hspace{3em}\textbf{Compute} $\mathcal{L}_{\text{CE}}$ from $z_{S}$ and $y$ \\
16 & \hspace{3em}\textbf{Compute} $\mathcal{L}_{\text{KL}}$ between $\text{softmax}(z_{T}/T)$ and $\text{softmax}(z_{S}/T)$ \\
17 & \hspace{3em}\textbf{if} $M = \text{Vanilla}$ \textbf{then} \\
18 & \hspace{4.5em}$\mathcal{L}_{\text{total}} \leftarrow \alpha_{\text{CE}}\mathcal{L}_{\text{CE}} + \beta_{\text{KL}}T^{2}\mathcal{L}_{\text{KL}}$ \\
19 & \hspace{3em}\textbf{else if} $M = \text{Attention Transfer}$ \textbf{then} \\
20 & \hspace{4.5em}\textbf{Extract} teacher and student feature maps from selected hooks \\
21 & \hspace{4.5em}\textbf{Compute} L2-normalized spatial attention maps \\
22 & \hspace{4.5em}\textbf{Interpolate} attention maps if spatial resolutions differ \\
23 & \hspace{4.5em}\textbf{Compute} $\mathcal{L}_{\text{AT}}$ between teacher and student attention maps \\
24 & \hspace{4.5em}$\mathcal{L}_{\text{total}} \leftarrow \alpha_{\text{CE}}\mathcal{L}_{\text{CE}} + \beta_{\text{KL}}T^{2}\mathcal{L}_{\text{KL}} + \gamma\mathcal{L}_{\text{AT}}$ \\
25 & \hspace{3em}\textbf{end if} \\
26 & \hspace{3em}\textbf{Update} $S$ using AdamW with $\mathcal{L}_{\text{total}}$ \\
27 & \hspace{1.5em}\textbf{end for} \\
28 & \hspace{1.5em}\textbf{Evaluate} $S$ on $D_{\text{val}}$ using macro-F1 \\
29 & \hspace{1.5em}\textbf{if} macro-F1 improves \textbf{then} \\
30 & \hspace{3em}$S_{\text{final}} \leftarrow S$ \\
31 & \hspace{3em}$\text{best\_score} \leftarrow \text{macro-F1}$ \\
32 & \hspace{1.5em}\textbf{end if} \\
33 & \hspace{1.5em}\textbf{Apply} scheduler and early stopping \\
34 & \textbf{end for} \\
35 & \textbf{return} $S_{\text{final}}$ \\
\end{tabular}\\[1pt]
\rule{\columnwidth}{0.9pt}
\caption{VKD and AT-KD training procedure.}
\label{alg:training}
\end{figure}

\subsection{Teacher and student models and training configuration}

The teacher is a SwinV2-Tiny transformer and the student a
MobileNetV3-Large network; Figs.~\ref{fig:1} and~\ref{fig:2} illustrate
both. Table~\ref{tab:teacher_params} lists the teacher training settings,
and Table~\ref{tab:train_params} lists the comparative VKD and AT-KD
settings, including the loss weights and temperature. We optimize with
AdamW under a cosine-annealing schedule with early stopping and select
the best checkpoint per fold by macro-F1.

\begin{table}[tbp]
\centering
\caption{Teacher training parameters.}
\label{tab:teacher_params}
\footnotesize
\renewcommand{\arraystretch}{1.2}
\setlength{\tabcolsep}{3pt}
\begin{tabular}{ll}
\toprule
\textbf{Parameter} & \textbf{Value} \\
\midrule
Model & SwinV2-Tiny \\
Pretrained Weights & ImageNet \\
Input Size & 256$\times$256 \\
Epochs & 50 \\
Optimizer & AdamW \\
Learning Rate & 1$\times$10\textsuperscript{-}\textsuperscript{4} \\
Weight Decay & 1$\times$10\textsuperscript{-}$^{2}$ \\
Scheduler & CosineAnnealingLR (Tmax = 50) \\
Loss Function & Cross-Entropy Loss \\
Early Stopping Patience & 50 epochs \\
Batch Size & 32 \\
\bottomrule
\end{tabular}
\end{table}
\begin{table}[tbp]
\centering
\caption{VKD and AT-KD training parameters.}
\label{tab:train_params}
\footnotesize
\renewcommand{\arraystretch}{1.2}
\setlength{\tabcolsep}{3pt}
\begin{tabular}{l c c c L{0.36\textwidth}}
\toprule
\textbf{Parameter} & \textbf{Symbol} & \textbf{VKD} & \textbf{AT-KD} & \textbf{Description} \\
\midrule
CE loss weight &  & 0.5 & 1.0 & Weight for hard-label cross-entropy supervision \\
KL loss weight &  & 0.5 & 0.5 & Weight for soft-target distillation via KL-divergence \\
AT loss scale & $\gamma$ & -- & 50.0 & Scaling coefficient for LAT to balance optimization dynamics \\
Temperature & T & 4.0 & 4.0 & Temperature parameter for soft-target generation \\
KL temperature scaling & T$^{2}$ & Applied & Applied & Scaling factor used to compensate for temperature-induced gradient scaling \\
Loss components &  &  &  & Main objective components used during student training \\
Feature hooks & -- & None & Teacher: layers.2/3; Student: blocks.4/5 & Intermediate layers used for attention map extraction and transfer \\
Optimizer & -- & AdamW & AdamW & Optimizer used for student model training \\
Learning rate & $\eta$ & 1 $\times$ 10\textsuperscript{-}\textsuperscript{4} & 1 $\times$ 10\textsuperscript{-}\textsuperscript{4} & Initial learning rate used during optimization \\
Weight decay & $\lambda$wd & 1 $\times$ 10\textsuperscript{-}$^{2}$ & 1 $\times$ 10\textsuperscript{-}$^{2}$ & Regularization coefficient used by AdamW \\
Scheduler & -- & CosineAnnealingLR & CosineAnnealingLR & Learning rate scheduler with Tmax = 100 \\
Early stopping & -- & 50 epochs & 50 epochs & Patience value used to prevent overfitting \\
\bottomrule
\end{tabular}
\end{table}

\section{Experiments}
\subsection{Datasets}

We use three public ALL datasets with deliberately different spatial
character. ALL-IDB-1 holds high-resolution peripheral blood-smear images
with multiple cells and broad context; ALL-IDB-2 holds small single-cell
crops; and the Subtype ALL set poses a harder three-class problem (L1,
L2, L3) \cite{ref44,ref45,ref46,ref48}. Fig.~\ref{fig:3} shows
representative samples, and Fig.~\ref{fig:eda} summarizes composition:
class counts (108, 260, and 300 images), class balance, and image-size
range. The two binary sets differ mainly in spatial layout, multi-cell
versus single-cell, whereas the Subtype set is both three-class and
mildly imbalanced (L2 is the minority at 63 of 300 images). Spanning
single-cell and multi-cell composition is what lets us isolate the
dependence of attention transfer on spatial complexity.

\begin{figure}[tbp]
\centering
\includegraphics[width=0.95\linewidth]{Figures/fig_eda.png}
\caption{Exploratory data analysis of the three benchmarks: per-class image counts, class-imbalance ratio, and image-size range. ALL-IDB-1 and ALL-IDB-2 are binary; the Subtype set is three-class with L2 as the minority.}
\label{fig:eda}
\end{figure}

\subsection{Preprocessing and augmentation}

We resize all images to a common input resolution and normalize them, and
apply standard label-preserving augmentations (flips, rotations, and
photometric jitter) during training to curb overfitting on the small
datasets. VKD and AT-KD share identical preprocessing, splits, teacher,
and student, so any measured difference traces to the distillation
objective alone.

\subsection{Evaluation metrics and protocol}

We measure classification performance with macro-averaged accuracy,
precision, recall, F1-score, and specificity, defined in
\eqref{eq:11}--\eqref{eq:15}; validation loss is the cross-entropy in
\eqref{eq:16}. We summarize discriminative performance with the area
under the ROC curve and average precision (\eqref{eq:17} and
\eqref{eq:18}), and assess interpretability with Grad-CAM (\eqref{eq:19})
and t-SNE. All results come from stratified five-fold cross-validation
with identical folds for every method.

\begin{equation}\label{eq:11}
Accuracy = \frac{TP + TN}{N}
\end{equation}
\begin{equation}\label{eq:12}
\text{Precision}  =  \frac{1}{C}\sum_{i = 1}^{C}\frac{\text{TP}_{i}}{\text{TP}_{i}  +  \text{FP}_{i}}
\end{equation}
\begin{equation}\label{eq:13}
\text{Recall}  =  \frac{1}{C}\sum_{i = 1}^{C}\frac{\text{TP}_{i}}{\text{TP}_{i}  +  \text{FN}_{i}}
\end{equation}
\begin{equation}\label{eq:14}
F_{1}  =  \frac{1}{C}\sum_{i = 1}^{C}\frac{2  \cdot  \text{Precision}_{i}  \cdot  \text{Recall}_{i}}{\text{Precision}_{i}  +  \text{Recall}_{i}}
\end{equation}
\begin{equation}\label{eq:15}
\text{Specificity} = \frac{1}{C}\sum_{i=1}^{C}\frac{\text{TN}_{i}}{\text{TN}_{i} + \text{FP}_{i}}
\end{equation}
\begin{equation}\label{eq:16}
\mathcal{L}_{\text{val}} = - \frac{1}{N}\sum_{i = 1}^{N}{\sum_{j = 1}^{C}y_{\text{ij}}}\log({\widehat{y}}_{\text{ij}})
\end{equation}
\begin{equation}\label{eq:17}
AUC = \int_{0}^{1}{TPR(FPR)\text{ }d(FPR)}
\end{equation}
\begin{equation}\label{eq:18}
\begin{aligned}
\text{AP} &= \sum_{k=1}^{n}\left( \text{Recall}_{k+1} - \text{Recall}_{k} \right) \cdot \text{Precision}_{k+1}
\end{aligned}
\end{equation}

\subsection{Statistical analysis}

To avoid over-reading fold-level differences, we apply four complementary
tests. The Friedman test assesses overall differences among methods; the
paired $t$-test (\eqref{eq:21}) evaluates fold-level differences between
AT-KD and VKD; McNemar's test (\eqref{eq:20}) compares paired prediction
errors; and Cohen's $d$ and Cliff's $\delta$ (\eqref{eq:22} and
\eqref{eq:23}) quantify effect size. Together they separate genuine
effects from cross-fold variance.

\begin{equation}\label{eq:19}
\begin{aligned}
\text{Heatmap} &= \text{ReLU}\!\left( \sum_{k}\alpha_{k} \cdot A^{k} \right), \\
\alpha_{k} &= \frac{1}{Z}\sum_{i,j}\frac{\partial y^{c}}{\partial A_{ij}^{k}}
\end{aligned}
\end{equation}
\begin{equation}\label{eq:20}
\chi^{2} = \frac{(\,|b-c|-1)^{2}}{b + c},H_{0}: \text{Error rates are identical}
\end{equation}
\begin{equation}\label{eq:21}
t = \frac{\bar{d}}{s_{d}/\sqrt{n}}
\end{equation}
\begin{equation}\label{eq:22}
d = \frac{\mu_{1} - \mu_{2}}{\sigma_{\text{pooled}}}
\end{equation}
\begin{equation}\label{eq:23}
\delta = \frac{P(X > Y) - P(X < Y)}{P(X > Y) + P(X < Y)}
\end{equation}

\subsection{The compact student matches the teacher, and AT-KD matches VKD}

Across the three datasets, the distilled MobileNetV3-Large student
reaches accuracy competitive with the SwinV2-Tiny teacher while using
only 4.21 million parameters, 6.6 times fewer parameters and 20.6 times
fewer FLOPs than the teacher (Table~\ref{tab:complexity}).
Table~\ref{tab:perf} reports per-dataset performance of VKD and AT-KD
under stratified five-fold cross-validation. On ALL-IDB-1, AT-KD reaches
96.36\% accuracy against 94.42\% for VKD; on ALL-IDB-2 the two stay
within one point (95.38\% for VKD, 94.62\% for AT-KD); and on the harder
three-class Subtype ALL task both reach 98.00\%. Per-class results follow
the same pattern (Table~\ref{tab:perf_class}), and the aggregated
confusion matrices show that the residual errors of the two methods fall
on the same class pairs (Fig.~\ref{fig:4}). The differences between AT-KD
and VKD are therefore small and favor neither method consistently across
datasets.

\begin{table}[tbp]
\centering
\caption{Teacher and student model complexity.}
\label{tab:complexity}
\footnotesize
\renewcommand{\arraystretch}{1.2}
\setlength{\tabcolsep}{3pt}
\begin{tabular}{l l c c c c}
\toprule
\textbf{Model} & \textbf{Role} & \textbf{Parameters (M)} & \textbf{Model Size (MB)} & \textbf{GFLOPs} & \textbf{GMACs} \\
\midrule
SwinV2-Tiny (window 8$\times$8) & Teacher & 27.58 & 107.06 & 11.92 & 5.96 \\
MobileNetV3-Large & Student & 4.21 & 16.14 & 0.58 & 0.29 \\
Compression Ratio &  & 6.6$\times$ & 6.6$\times$ & 20.6$\times$ & 20.6$\times$ \\
\bottomrule
\end{tabular}
\end{table}
\begin{table}[tbp]
\centering
\caption{Classification performance of VKD and AT-KD (mean $\pm$ std, 5-fold CV).}
\label{tab:perf}
\footnotesize
\renewcommand{\arraystretch}{1.2}
\setlength{\tabcolsep}{3pt}
\begin{tabular}{l l c c c c c}
\toprule
\textbf{Dataset} & \textbf{Method} & \textbf{Accuracy (\%)} & \textbf{Precision (\%)} & \textbf{Recall (\%)} & \textbf{F1-score (\%)} & \textbf{Val Loss} \\
\midrule
ALL-IDB-1 & VKD & 94.42 $\pm$ 5.99 & 94.68 $\pm$ 5.67 & 94.52 $\pm$ 6.13 & 94.35 $\pm$ 6.08 & 0.381 $\pm$ 0.415 \\
 & AT-KD & 96.36 $\pm$ 4.98 & 96.50 $\pm$ 4.80 & 96.50 $\pm$ 4.80 & 96.35 $\pm$ 5.00 & 0.163 $\pm$ 0.153 \\
ALL-IDB-2 & VKD & 95.38 $\pm$ 2.19 & 95.47 $\pm$ 2.23 & 95.38 $\pm$ 2.19 & 95.38 $\pm$ 2.19 & 0.252 $\pm$ 0.098 \\
 & AT-KD & 94.62 $\pm$ 4.98 & 94.74 $\pm$ 4.93 & 94.62 $\pm$ 4.98 & 94.61 $\pm$ 4.98 & 0.224 $\pm$ 0.139 \\
Subtype ALL & VKD & 98.00 $\pm$ 1.39 & 97.55 $\pm$ 1.93 & 97.90 $\pm$ 1.50 & 97.65 $\pm$ 1.63 & 0.115 $\pm$ 0.079 \\
 & AT-KD & 98.00 $\pm$ 1.83 & 97.76 $\pm$ 2.34 & 97.69 $\pm$ 1.90 & 97.67 $\pm$ 2.06 & 0.082 $\pm$ 0.037 \\
\bottomrule
\end{tabular}
\end{table}
\begin{table}[tbp]
\centering
\caption{Per-class classification performance (mean $\pm$ std, 5-fold CV).}
\label{tab:perf_class}
\footnotesize
\renewcommand{\arraystretch}{1.2}
\setlength{\tabcolsep}{3pt}
\begin{tabular}{l l l c c c c}
\toprule
\textbf{Dataset} & \textbf{Class} & \textbf{Method} & \textbf{Precision (\%)} & \textbf{Recall (\%)} & \textbf{F1-score (\%)} & \textbf{Support} \\
\midrule
ALL-IDB-1 & Cancer & VKD & 92.44 $\pm$ 6.27 & 96.00 $\pm$ 8.94 & 93.97 $\pm$ 6.24 & 9--10 \\
 &  & AT-KD & 94.67 $\pm$ 6.29 & 98.00 $\pm$ 4.47 & 96.18 $\pm$ 4.58 & 9--10 \\
 & Non-cancer & VKD & 96.92 $\pm$ 5.84 & 93.03 $\pm$ 7.63 & 94.73 $\pm$ 6.04 & 11--12 \\
 &  & AT-KD & 96.33 $\pm$ 4.04 & 95.00 $\pm$ 6.67 & 96.53 $\pm$ 5.41 & 11--12 \\
ALL-IDB-2 & Blast & VKD & 94.78 $\pm$ 2.88 & 96.15 $\pm$ 2.59 & 95.43 $\pm$ 1.94 & 26 \\
 &  & AT-KD & 94.92 $\pm$ 6.12 & 95.38 $\pm$ 3.37 & 94.70 $\pm$ 4.34 & 26 \\
 & Non-blast & VKD & 96.15 $\pm$ 2.90 & 94.62 $\pm$ 3.37 & 95.34 $\pm$ 2.45 & 26 \\
 &  & AT-KD & 94.56 $\pm$ 3.74 & 93.85 $\pm$ 6.63 & 94.52 $\pm$ 5.60 & 26 \\
Subtype ALL & L1 & VKD & 98.46 $\pm$ 1.98 & 96.92 $\pm$ 3.05 & 97.65 $\pm$ 1.59 & 25--26 \\
 &  & AT-KD & 98.43 $\pm$ 2.01 & 97.69 $\pm$ 3.37 & 97.85 $\pm$ 1.97 & 25--26 \\
 & L2 & VKD & 94.18 $\pm$ 5.32 & 96.77 $\pm$ 3.85 & 95.27 $\pm$ 3.05 & 12--13 \\
 &  & AT-KD & 95.60 $\pm$ 5.60 & 94.92 $\pm$ 3.85 & 95.38 $\pm$ 4.19 & 12--13 \\
 & L3 & VKD & 100.00 $\pm$ 0.00 & 100.00 $\pm$ 0.00 & 100.00 $\pm$ 0.00 & 21--22 \\
 &  & AT-KD & 100.00 $\pm$ 0.00 & 100.00 $\pm$ 0.00 & 100.00 $\pm$ 0.00 & 21--22 \\
\bottomrule
\end{tabular}
\end{table}
\begin{figure}[tbp]
\centering
\includegraphics[width=0.8\linewidth]{Figures/fig4.png}
\caption{Aggregated confusion matrices for VKD and AT-KD on each dataset.}
\label{fig:4}
\end{figure}

\subsection{Attention transfer lowers validation loss and accelerates convergence}

The clearest and most consistent effect of attention transfer is on
optimization, not accuracy. AT-KD reaches a markedly lower validation
loss than VKD on the spatially complex datasets: on ALL-IDB-1 it falls
from 0.381 to 0.163, and on Subtype ALL from 0.115 to 0.082
(Table~\ref{tab:perf}). The learning curves confirm that AT-KD reaches a
lower, more stable validation loss with smaller fold-to-fold variance
(Fig.~\ref{fig:5}), and the training-time analysis shows it does so in
comparable or fewer epochs and at no inference cost, since the
attention-transfer terms run only during training
(Table~\ref{tab:traintime}). Discriminative performance, summarized by
the ROC and precision--recall curves, is high for both methods and
slightly better for AT-KD on the high-resolution datasets
(Figs.~\ref{fig:6} and~\ref{fig:7}).

\begin{table}[tbp]
\centering
\caption{Training time of VKD and AT-KD (5-fold CV).}
\label{tab:traintime}
\footnotesize
\renewcommand{\arraystretch}{1.2}
\setlength{\tabcolsep}{3pt}
\begin{tabular}{l l c c c}
\toprule
\textbf{Dataset} & \textbf{Method} & \textbf{Mean Total (s)} & \textbf{Mean Epochs} & \textbf{Time/Epoch (s)} \\
\midrule
ALL-IDB-1 & VKD & 267.28 & 78.2 & 3.43 \\
 & AT-KD & 234.14 & 67.6 & 3.47 \\
ALL-IDB-2 & VKD & 142.94 & 72.0 & 1.98 \\
 & AT-KD & 133.21 & 66.0 & 2.02 \\
Subtype ALL & VKD & 176.75 & 67.6 & 2.62 \\
 & AT-KD & 180.98 & 67.6 & 2.68 \\
\bottomrule
\end{tabular}
\end{table}
\begin{figure}[tbp]
\centering
\includegraphics[width=0.85\linewidth]{Figures/fig5.png}
\caption{Training and validation learning curves for VKD and AT-KD.}
\label{fig:5}
\end{figure}
\begin{figure}[tbp]
\centering
\includegraphics[width=0.7\linewidth]{Figures/fig6.png}
\caption{ROC curves of VKD and AT-KD on the three datasets.}
\label{fig:6}
\end{figure}
\begin{figure}[tbp]
\centering
\includegraphics[width=0.7\linewidth]{Figures/fig7.png}
\caption{Precision--Recall curves of VKD and AT-KD on the three datasets.}
\label{fig:7}
\end{figure}

\subsection{Statistical analysis confirms accuracy equivalence}

We tested the accuracy difference between AT-KD and VKD with four
complementary tests. The Friedman test detects an overall difference
among methods on the two binary datasets but not on Subtype ALL
(Table~\ref{tab:friedman}). The direct paired comparisons between AT-KD
and VKD, however, are not significant for any metric on any dataset:
paired $t$-test accuracy $p$-values are 0.569, 0.648, and 1.000, and
every McNemar test is likewise non-significant (Table~\ref{tab:ttest},
Table~\ref{tab:mcnemar}, Table~\ref{tab:signif}). Effect sizes are
correspondingly small to negligible (Cohen's $d$ and Cliff's $\delta$;
Table~\ref{tab:effect}) and sit within or near the negligible band on
every dataset (Fig.~\ref{fig:9}). The accuracy of AT-KD and VKD is
therefore statistically equivalent, and we read this parity as evidence
that both already recover the discriminative signal available at this
model scale.

\begin{table}[tbp]
\centering
\caption{Friedman test on the three datasets.}
\label{tab:friedman}
\footnotesize
\renewcommand{\arraystretch}{1.2}
\setlength{\tabcolsep}{3pt}
\begin{tabular}{l l c c c}
\toprule
\textbf{Dataset} & \textbf{Metric} & $\chi^{2}$ & \textbf{p-value} & \textbf{Significant (}$\alpha$\textbf{=0.05)?} \\
\midrule
ALL-IDB-1 & Accuracy & 11.857 & 0.0368 & Yes \\
 & F1-score & 11.857 & 0.0368 & Yes \\
ALL-IDB-2 & Accuracy & 16.468 & 0.0056 & Yes \\
 & F1-score & 16.468 & 0.0056 & Yes \\
Subtype-ALL & Accuracy & 3.222 & 0.6658 & No \\
 & F1-score & 5.208 & 0.3910 & No \\
\bottomrule
\end{tabular}
\end{table}
\begin{table}[tbp]
\centering
\caption{Paired t-test between VKD and AT-KD.}
\label{tab:ttest}
\footnotesize
\renewcommand{\arraystretch}{1.2}
\setlength{\tabcolsep}{3pt}
\begin{tabular}{l l c c c c c c}
\toprule
\textbf{Dataset} & \textbf{Metric} & \textbf{Mean Diff} & \textbf{Std Diff} & \textbf{t-stat} & \textbf{p-value} & \textbf{Sig?} & \textbf{CI 95\%} \\
\midrule
ALL-IDB-1 & Accuracy & -0.0195 & 0.0703 & -0.620 & 0.5691 & No & [-0.107, 0.068] \\
 & Precision & -0.0182 & 0.0682 & -0.595 & 0.5837 & No & [-0.103, 0.067] \\
 & Recall & -0.0198 & 0.0704 & -0.630 & 0.5627 & No & [-0.107, 0.068] \\
 & F1-score & -0.0200 & 0.0710 & -0.630 & 0.5631 & No & [-0.108, 0.068] \\
ALL-IDB-2 & Accuracy & +0.0077 & 0.0349 & +0.492 & 0.6483 & No & [-0.036, 0.051] \\
 & Precision & +0.0072 & 0.0335 & +0.483 & 0.6543 & No & [-0.034, 0.049] \\
 & Recall & +0.0077 & 0.0349 & +0.492 & 0.6483 & No & [-0.036, 0.051] \\
 & F1-score & +0.0077 & 0.0350 & +0.494 & 0.6473 & No & [-0.036, 0.051] \\
Subtype-ALL & Accuracy & 0.0000 & 0.0236 & 0.000 & 1.0000 & No & [-0.029, 0.029] \\
 & Precision & -0.0021 & 0.0267 & -0.176 & 0.8685 & No & [-0.035, 0.031] \\
 & Recall & +0.0020 & 0.0284 & +0.160 & 0.8806 & No & [-0.033, 0.037] \\
 & F1-score & -0.0003 & 0.0273 & -0.022 & 0.9833 & No & [-0.034, 0.034] \\
\bottomrule
\end{tabular}
\end{table}
\begin{table}[tbp]
\centering
\caption{McNemar's test between VKD and AT-KD.}
\label{tab:mcnemar}
\footnotesize
\renewcommand{\arraystretch}{1.2}
\setlength{\tabcolsep}{3pt}
\begin{tabular}{l c c c c c c c L{0.16\textwidth}}
\toprule
\textbf{Dataset} & \textbf{Samples} & \textbf{Both corr.} & \textbf{Both wrong} & \textbf{VKD only ($b$)} & \textbf{AT-KD only ($c$)} & \textbf{p-value} & \textbf{Sig?} & \textbf{Interpretation} \\
\midrule
ALL-IDB-1 & 108 & 98 & 0 & 4 & 6 & 0.7539 & No & Similar prediction pattern \\
ALL-IDB-2 & 260 & 237 & 3 & 11 & 9 & 0.8238 & No & Similar prediction pattern \\
Subtype-ALL & 300 & 290 & 2 & 4 & 4 & 1.0000 & No & Nearly identical observed prediction pattern \\
\bottomrule
\end{tabular}
\end{table}
\begin{table}[tbp]
\centering
\caption{Effect sizes between VKD and AT-KD.}
\label{tab:effect}
\footnotesize
\renewcommand{\arraystretch}{1.2}
\setlength{\tabcolsep}{3pt}
\begin{tabular}{l c c c c c c c c}
\toprule
\textbf{Dataset} & \textbf{Metric} & \textbf{Cohen's $d$} & \textbf{Interp.} & \textbf{Cliff's $\delta$} & \textbf{Interp.} & \textbf{Wins} & \textbf{Ties} & \textbf{Losses} \\
\midrule
ALL-IDB-1 & Accuracy & -0.354 & Small & -0.28 & Small & 1 & 1 & 3 \\
 & F1-score & -0.359 & Small & -0.28 & Small & 1 & 1 & 3 \\
ALL-IDB-2 & Accuracy & +0.200 & Negligible & 0.00 & Negligible & 3 & 0 & 2 \\
 & F1-score & +0.201 & Small & +0.04 & Negligible & 3 & 0 & 2 \\
Subtype-ALL & Accuracy & 0.000 & Negligible & -0.08 & Negligible & 2 & 0 & 3 \\
 & F1-score & -0.015 & Negligible & -0.08 & Negligible & 2 & 0 & 3 \\
\bottomrule
\end{tabular}
\end{table}
\begin{table}[tbp]
\centering
\caption{Summary of cross-method significance.}
\label{tab:signif}
\footnotesize
\renewcommand{\arraystretch}{1.2}
\setlength{\tabcolsep}{3pt}
\begin{tabular}{L{0.26\textwidth} L{0.22\textwidth} L{0.22\textwidth} L{0.20\textwidth}}
\toprule
\textbf{Aspect} & \textbf{ALL-IDB-1} & \textbf{ALL-IDB-2} & \textbf{Subtype-ALL} \\
\midrule
Friedman (global) & Significant* (p = 0.037) & Significant* (p = 0.006) & Not significant \\
Paired t-Test: Accuracy & p = 0.569 (No) & p = 0.648 (No) & p = 1.000 (No) \\
Paired t-Test: Precision & p = 0.584 (No) & p = 0.654 (No) & p = 0.869 (No) \\
Paired t-Test: Recall & p = 0.563 (No) & p = 0.648 (No) & p = 0.881 (No) \\
Paired t-Test: F1-score & p = 0.563 (No) & p = 0.647 (No) & p = 0.983 (No) \\
McNemar & p = 0.754 (No) & p = 0.824 (No) & p = 1.000 (No) \\
Cohen's d (Accuracy) & -0.354 (small) & +0.200 (negligible) & 0.000 (negligible) \\
Cohen's d (F1-score) & -0.359 (small) & +0.201 (small) & -0.015 (negligible) \\
Cliff's $\delta$ (Accuracy) & -0.28 (small) & 0.00 (negligible) & -0.08 (negligible) \\
Cliff's $\delta$ (F1-score) & -0.28 (small) & +0.04 (negligible) & -0.08 (negligible) \\
Direction of Advantage & AT-KD slightly superior & VKD slightly superior & Comparable \\
\bottomrule
\end{tabular}
\end{table}
\begin{figure}[tbp]
\centering
\includegraphics[width=0.6\linewidth]{Figures/fig9.png}
\caption{Effect sizes (Cohen's $d$ and Cliff's $\delta$) for the AT-KD versus VKD difference.}
\label{fig:9}
\end{figure}

\subsection{The benefit of attention transfer depends on image spatial complexity}

The advantage of attention transfer is not uniform across datasets. Its
gains in validation loss and focused activation are clear on ALL-IDB-1
and Subtype ALL, which hold high-resolution images with multiple cells
and broader context, and negligible on ALL-IDB-2, which holds small
single-cell crops. Because we resize all inputs to a common resolution
before training, this pattern reflects the spatial complexity and cell
composition of the source images rather than the network input resolution
itself: when an image holds little structure beyond a single centred
cell, there is little spatial attention for the teacher to transfer.

\subsection{Attention transfer improves interpretability}

Beyond the quantitative metrics, attention transfer changes where and how
the student looks. Grad-CAM maps show that the AT-KD student concentrates
on morphologically relevant regions, the nuclei, cytoplasm, and
blast-cell clusters, whereas the VKD student produces more diffuse
activation (Fig.~\ref{fig:10}). The t-SNE projections of the
penultimate-layer features show tighter intra-class clusters and clearer
inter-class separation for AT-KD (Fig.~\ref{fig:11}). The complexity
comparison in Fig.~\ref{fig:8} sets these gains against the large
reduction in model size. Together with the lower validation loss, these
visualizations indicate that attention transfer yields a student whose
internal reasoning aligns better with diagnostic morphology, which
matters for clinical trust even when accuracy is unchanged.

\begin{figure}[tbp]
\centering
\includegraphics[width=0.55\linewidth]{Figures/fig8.png}
\caption{Model complexity of the SwinV2-Tiny teacher and the MobileNetV3-Large student.}
\label{fig:8}
\end{figure}
\begin{figure}[tbp]
\centering
\includegraphics[width=0.7\linewidth]{Figures/fig10.png}
\caption{Grad-CAM attention maps for VKD and AT-KD on representative samples.}
\label{fig:10}
\end{figure}
\begin{figure}[tbp]
\centering
\includegraphics[width=0.6\linewidth]{Figures/fig11.png}
\caption{t-SNE projections of the student's penultimate-layer features.}
\label{fig:11}
\end{figure}

\subsection{Comparison with reported methods}

Table~\ref{tab:sota} sets the distilled student against previously
reported methods on the same datasets, and Table~\ref{tab:method_comp}
summarizes the methodological choices of representative studies. The
compact student stays competitive with much larger reported pipelines
while remaining far smaller, and, unlike most prior work, our evaluation
reports both the positive and the limited effects of attention transfer
under a uniform cross-validation and statistical protocol rather than a
single optimistic split.

\begin{table}[tbp]
\centering
\caption{Comparison with state-of-the-art baselines.}
\label{tab:sota}
\footnotesize
\renewcommand{\arraystretch}{1.2}
\setlength{\tabcolsep}{3pt}
\begin{tabular}{L{0.22\textwidth} L{0.18\textwidth} c c c c c}
\toprule
\textbf{Method} & \textbf{Dataset} & \textbf{Accuracy} & \textbf{Precision} & \textbf{Sensitivity} & \textbf{Specificity} & \textbf{F1-score} \\
\midrule
Attention Transfer KD (Ours) & Subtype ALL (5-Fold CV) & 98.00\% & 97.76\% & 97.69\% & 99.00\% & 97.67\% \\
VKD (Ours) & Subtype ALL (5-Fold CV) & 98.00\% & 97.55\% & 97.90\% & 99.05\% & 97.65\% \\
Attention Transfer KD (Ours) & ALL-IDB-1 (5-Fold CV) & 96.36\% & 96.50\% & 96.50\% & 95.00\% & 96.35\% \\
VKD (Ours) & ALL-IDB-2 (5-Fold CV) & 95.38\% & 95.47\% & 95.38\% & 94.62\% & 95.38\% \\
Attention Transfer KD (Ours) & ALL-IDB-2 (5-Fold CV) & 94.62\% & 94.74\% & 94.62\% & 94.62\% & 94.61\% \\
VKD (Ours) & ALL-IDB-1 (5-Fold CV) & 94.42\% & 94.68\% & 94.52\% & 93.03\% & 94.35\% \\
Grey-scaling Normalization + SVM \cite{ref52} & ALL-IDB-2 & 92.80\% & - & - & - & - \\
LeuFeatx + XGB (5-fold CV) \cite{ref47} & ALL-IDB-2 & 92.67\% & 91.62\% & 94.96\% & 89.87\% & 93.20\% \\
K-NN (7 selected features) \cite{ref53} & ALL-IDB-2 (70\% training, 30\% testing) & 92.41\% & 93.00\% & 92.86\% & 91.89\% & - \\
VGG16-OSLNet \cite{ref54} & ALL-IDB-2 (70\% training, 30\% testing) & 92.31\% & 91.88\% & 92.82\% & 91.79\% & 92.35\% \\
ResNet50-OSLNet \cite{ref54} & ALL-IDB-2 (10-Fold CV) & 92.31\% & 93.01\% & 91.54\% & 93.08\% & 92.27\% \\
LeuFeatx + SVM (5-fold CV) \cite{ref47} & ALL-IDB-2 & 92.30\% & 90.63\% & 96.67\% & 86.36\% & 93.55\% \\
ResNet50 (Pre-trained) \cite{ref55} & ALL-IDB-2 (80\% Training, 20\% Testing) & 92.30\% & 92.55\% & 92.25\% & 91.80\% & - \\
Xception (Transfer Learning - Centralized) \cite{ref56} & PBS Dataset (Taleqani Hospital, Iran) - (70\% train, 20\% test, 10\% val) & 91.80\% & 92.40\% & 91.10\% & - & 91.75\% \\
DenseNet121 (Transfer Learning - Centralized) \cite{ref56} & PBS Dataset (Taleqani Hospital, Iran) - (70\% train, 20\% test, 10\% val) & 91.23\% & 91.80\% & 90.90\% & - & 91.35\% \\
\bottomrule
\end{tabular}
\end{table}
\begin{table}[tbp]
\centering
\caption{Methodological comparison of representative state-of-the-art studies.}
\label{tab:method_comp}
\footnotesize
\renewcommand{\arraystretch}{1.2}
\setlength{\tabcolsep}{3pt}
\begin{tabular}{L{0.28\textwidth} L{0.22\textwidth} L{0.40\textwidth}}
\toprule
\textbf{Method} & \textbf{Dataset(s) Used} & \textbf{Contribution} \\
\midrule
Fine-tuned VGG16 feature extractor + SVM / RF classifier \cite{ref47} & AML Morph., PBC Barcelona, ALL-IDB-2 & Cross-dataset feature extractor with interpretable filter visualisation for cell-level reasoning. \\
Enhanced edge detection for touching-cell segmentation + subtype classifier \cite{ref48} & Subtype ALL (L1, L2, L3) & Addresses overlapping-cell segmentation, a failure mode often unaddressed in conventional pipelines. \\
Hybrid MobileNetV2 + ResNet18 with probability-based weight fusion \cite{ref51} & ALL-IDB-1, ALL-IDB-2 & Probability-based weight factor for fusing heterogeneous CNN backbones in ALL classification. \\
Classical CAD pipeline: morphological + texture features to SVM / ANN \cite{ref52} & ALL-IDB & Foundational soft-computing CAD reference for automated ALL diagnosis. \\
Handcrafted features + ACO feature selection + NB / SVM / K-NN \cite{ref53} & ALL-IDB-2 (260 imgs) & Establishes optimal feature selection as a competitive alternative to deep models on small datasets. \\
ResNet18 + Orthogonal SoftMax Layer (OSL) classifier \cite{ref54} & ALL-IDB-1/2, C-NMC 2019, ASH & Orthogonal-weight classifier mitigates overfitting on small medical datasets, applicable to both ALL and AML. \\
Active-contour + 3-CNN feature fusion + PCA + RF / XGBoost \cite{ref55} & C-NMC 2019, ALL-IDB-2 & Serial multi-CNN feature fusion with dimensionality reduction for compact representation learning. \\
Federated Learning + differential privacy across InceptionV3 / DenseNet121 / Xception \cite{ref56} & PBS Taleqani Tehran (3,242 imgs) & multi-class federated-learning framework for ALL subtypes that preserves patient privacy under non-IID conditions. \\
Attention Transfer Knowledge Distillation (AT-KD): SwinV2-Tiny teacher to MobileNetV3-Large student transferring softened logits AND L2-normalized spatial attention maps from selected stages, Proposed Model & Subtype ALL, ALL-IDB-1, ALL-IDB-2 (stratified 5-fold CV) & Unifies accuracy, model compression, optimisation stability, and intrinsic spatial interpretability within a single distillation framework, transferring teacher attention without pixel-level annotation and incurring no inference-time overhead. To the best of our knowledge, the only approach in the surveyed literature that jointly addresses all four objectives. \\
\bottomrule
\end{tabular}
\end{table}

\section{Discussion}

The central result is that, for lightweight ALL classification, accuracy
cannot settle the choice between attention-transfer and vanilla
distillation. Across three datasets, four statistical tests, and both
binary and three-class settings, AT-KD and VKD are statistically
indistinguishable, with effect sizes that never leave the
small-to-negligible range. We read this parity not as a failure of
attention transfer but as a property of the regime: at the scale of a
4.21-million-parameter student trained on these benchmarks, the
discriminative signal already saturates, so any well-posed distillation
objective recovers it. An evaluation that stops at accuracy cannot
separate the two methods, and a study that reports only accuracy gains
risks crediting attention transfer with an effect the data do not
support.

Where the two methods diverge is in optimization. AT-KD reaches a lower
validation loss and converges in comparable or fewer epochs, most clearly
on the spatially complex datasets, where the ALL-IDB-1 validation loss
falls by more than half relative to vanilla distillation. A plausible
mechanism is that aligning the student's intermediate spatial attention
with the teacher's adds a geometry-aware constraint that logit matching
alone cannot supply; this constraint narrows the solutions the student
can adopt and steers it toward representations organized around
diagnostic structure, smoothing the loss surface and reducing
fold-to-fold variance. Because this auxiliary signal acts only during
training, the deployed student is identical to a vanilla-distilled one, so
the optimization benefit comes at zero inference cost. The same mechanism
explains why the benefit does not surface as accuracy: on saturated data,
a smoother path to the same optimum changes how the model is reached, not
where it ends up.

The dependence of this benefit on the input is the most actionable
finding. Gains in validation loss and focused activation are pronounced on
ALL-IDB-1 and the three-class Subtype ALL set, both of which hold
high-resolution fields with multiple cells and surrounding context, and
they essentially vanish on ALL-IDB-2, whose images are tightly cropped
single cells. Because we resize every input to a common resolution before
training, this pattern reflects the spatial complexity and cell
composition of the source image rather than its pixel count: when an
image holds little structure beyond one centred cell, there is almost no
spatial attention map for the teacher to transfer, and the auxiliary loss
degenerates toward a constant. This yields a concrete decision rule
rather than a vague recommendation: attention transfer is worth its
training cost when the imagery carries multi-cell spatial context and can
be omitted when it does not.

The interpretability evidence sharpens this picture rather than
decorating it. The Grad-CAM maps show that AT-KD concentrates the
student's response on nuclei, cytoplasm, and blast-cell clusters, whereas
vanilla distillation produces more diffuse activation, and the t-SNE
projections show tighter intra-class clusters and cleaner inter-class
margins for AT-KD. For a deployment-oriented model these properties
matter independently of accuracy, because a clinician's trust depends on
whether a model attends to the features a hematologist would use, not
only on whether its label is correct. We stay deliberately cautious here:
visual alignment with morphology is necessary but not sufficient evidence
of clinically valid reasoning, and the present results establish that
attention transfer changes where the model looks, not yet that the change
improves diagnostic safety.

Several limitations bound these conclusions. The three datasets are
widely used public benchmarks of modest size, and their high absolute
accuracies leave little headroom in which an accuracy difference could
emerge even if one existed; a fair test of whether attention transfer can
convert its optimization advantage into accuracy needs larger and harder
data, such as the C-NMC challenge corpus, and external cohorts that vary
staining, scanner, and site, including privacy-preserving and federated
deployments \cite{ref50}. The analysis also fixes a single
teacher--student pair and trains the student from scratch without
large-scale pretraining, so teacher variance and the interaction between
attention transfer and pretrained initialization stay unquantified.
Calibration, although favorable for the distilled student, was measured
rather than directly optimized, and predictive uncertainty, central to
safe screening, was not estimated. The statistical protocol controls
fold-level variance well but cannot rule out small effects that only a
larger sample would reveal.

These limitations define a programme rather than a caveat. The most
informative next step is to test the optimization advantage on data with
genuine headroom, where a lower validation loss and faster convergence
could plausibly translate into measurable accuracy or calibration gains.
Pairing attention transfer with calibration-aware training and
uncertainty estimation, and with saliency-guided augmentation that
supplies additional spatial supervision \cite{ref49}, is a natural way to
push the conditional benefit we observe toward a consistent one. More
broadly, our results argue for evaluating distillation on a multi-axis
basis spanning optimization, calibration, interpretability, and
efficiency rather than on accuracy alone, because that is where the
methodological differences between distillation schemes actually reside.

\section{Conclusion}
We presented a controlled, multi-lens study of attention transfer for
lightweight ALL classification, and three findings stand. A
4.21-million-parameter student matches its transformer teacher at 6.6
times fewer parameters and 20.6 times fewer FLOPs. AT-KD and VKD are
statistically equivalent in accuracy, which four complementary tests show
is the expected outcome on saturated benchmarks rather than a shortcoming
of attention transfer. And the genuine value of attention transfer is an
optimization one, lower validation loss, faster convergence, and tighter,
more morphology-aligned representations, conditional on multi-cell
spatial complexity and obtained at zero inference cost. Together these
yield a concrete rule for when attention transfer earns its training cost
and a reproducible protocol for evaluating distillation across
optimization, interpretability, and efficiency rather than accuracy
alone. Future work should extend the study to larger external corpora
such as the C-NMC challenge data, fold teacher variance and pretraining
into the protocol, and pair attention transfer with calibration-aware
training and uncertainty estimation \cite{ref49} to test whether its
optimization advantage can become a measurable accuracy or calibration
gain.

\section*{Data availability}
The three datasets analysed in this study are publicly available:
ALL-IDB-1 and ALL-IDB-2 from the ALL-IDB database, and the Subtype ALL
dataset from its original release. The code and per-fold analysis
artifacts that support the findings are available from the corresponding
author on reasonable request.

\section*{Acknowledgements}
This work was supported by the Indonesian Endowment Fund for Education
(LPDP) on behalf of the Indonesian Ministry of Higher Education, Science
and Technology, under the EQUITY Program (Contract No.
4299/B3/DT.03.08/2025 and No. 3029/PKS/ITS/2025). The authors used a
large language model (Claude, Anthropic) solely to assist with language
editing; all scientific content is the authors' own.

\section*{CRediT authorship contribution statement}
\textbf{Tanzilal Mustaqim:} Conceptualization, Methodology, Software,
Investigation, Writing -- original draft. \textbf{Chastine Fatichah:}
Supervision, Writing -- review \& editing. \textbf{Nanik Suciati:}
Supervision, Writing -- review \& editing. \textbf{Takashi Obi:}
Methodology, Formal analysis. \textbf{Joong-Sun Lee:} Methodology, Formal
analysis. All authors reviewed and approved the manuscript.

\section*{Declaration of competing interest}
The authors declare no competing interests.

\bibliographystyle{elsarticle-num}
\bibliography{references}

\end{document}
